\documentclass[UTF8, a4paper, 12pt]{article}

\usepackage{fontspec}
\setmainfont{Times New Roman}
\usepackage{xeCJK}
\setCJKmainfont{SimSun}

\usepackage{geometry}
\geometry{margin=1in, top=1.2in, bottom=1.2in}

\usepackage{amsmath, amssymb, amsfonts}
\usepackage{graphicx}
\usepackage{booktabs}
\usepackage{dcolumn}
\usepackage{siunitx}
\usepackage{hyperref}
\usepackage[square,super]{natbib}

\hypersetup{
    colorlinks=false,
    linkcolor=black,
    citecolor=black,
    urlcolor=black,
    filecolor=black,
}

\title{基于多项式残差克里金与Matern核高斯过程回归的PM$_{2.5}$ CMAQ数据融合方法}
\author{研究团队}
\date{\today}

\setlength{\parindent}{2em}
\setlength{\parskip}{0pt}

\begin{document}

\maketitle

\begin{abstract}
PM$_{2.5}$浓度的高精度估算对公共健康和大气污染防控具有重要意义。化学输送模型（CMAQ）能够提供高时空分辨率的模拟数据，但存在系统性偏差。本研究提出一种基于多项式残差克里金（Polynomial Residual Kriging）与Matern核高斯过程回归的PM$_{2.5}$数据融合方法（AdvancedRK）。该方法首先通过二阶多项式建立CMAQ模拟值与监测站观测值之间的非线性偏差校正模型，随后利用高斯过程回归（GPR）与Matern核函数（$\nu=1.5$）对残差进行空间插值，从而实现两类数据源的优势互补。针对传统RBF核函数在刻画PM$_{2.5}$空间相关性上的局限，本研究引入Matern核函数以更准确地描述污染物的空间扩散特征。基于中国某地区1528个地面监测站的多阶段十折交叉验证实验表明，AdvancedRK方法在预实验阶段（5天）的决定系数$R^2$达到0.9015，在1月整月验证（31天）中达到0.9143，在7月整月验证（31天）中达到0.8501，在12月整月验证（31天）中达到0.9104，全部四个阶段均通过创新判定，综合性能显著优于基准VNA方法。

\textbf{关键词：} PM$_{2.5}$；CMAQ；数据融合；多项式回归；高斯过程回归；Matern核函数；残差克里金
\end{abstract}

\section{引言}

大气颗粒物污染是全球面临的重大环境问题之一。PM$_{2.5}$（空气动力学直径不超过2.5微米的颗粒物）因其粒径小、比表面积大，可深入人体肺部甚至进入血液循环系统，与心血管疾病、呼吸系统疾病等健康问题密切相关~\cite{brooks2015review}。世界卫生组织（WHO）数据显示，全球每年约有700万人因空气污染过早死亡，其中PM$_{2.5}$是最主要的健康风险因子之一。实时准确的PM$_{2.5}$浓度监测对于疾病风险预警、污染防控决策和公众健康保护具有重要意义。

目前PM$_{2.5}$监测主要依赖两类数据源：地面监测站观测数据和化学输送模型（CTM）模拟数据~\cite{zheng2020kriging}。地面监测站提供精确的观测值，但站点分布稀疏，难以覆盖广大区域；CMAQ模型可生成高时空分辨率的连续格点数据，但存在系统性偏差~\cite{foley2010CMAQ}。数据融合技术通过结合两类数据源的优点，可生成既具有观测精度又覆盖空间全区域的高质量PM$_{2.5}$浓度数据~\cite{reich2014ensemble}。

传统的空间插值方法如最近邻法（VNA）直接应用于监测数据~\cite{singh2019enhanced}，忽略了CMAQ模型提供的空间结构信息。变分插值方法（VNA/eVNA/aVNA）通过引入协变量（如CMAQ模拟值）作为辅助变量~\cite{liu2019spatial}，在一定程度上改善了融合效果。然而，这些方法仍存在以下局限：（1）假设监测值与协变量之间存在线性关系，忽略了非线性依赖；（2）难以准确刻画残差的空间相关性结构~\cite{cooley2007kriging}。

为解决上述问题，本研究提出基于多项式残差克里金（Polynomial Residual Kriging）的PM$_{2.5}$ CMAQ融合方法。该方法的核心创新在于：（1）采用二阶多项式建模CMAQ与监测值之间的非线性偏差关系；（2）利用高斯过程回归（GPR）克里金对残差进行空间插值。进一步地，为更准确地描述PM$_{2.5}$的空间扩散特征，本研究引入Matern核函数（$\nu=1.5$）替代传统的RBF核函数，提出AdvancedRK方法~\cite{rasmussen2006gaussian}。

\section{数据与方法}

\subsection{研究区域与数据}

本研究使用中国某地区1528个地面监测站的PM$_{2.5}$日均浓度数据，以及对应时空分辨率的CMAQ模型模拟数据~\cite{foley2010CMAQ}。数据时间跨度为2020年全年。监测站分布涵盖城区、郊区和农村等不同下垫面类型，具有较好的空间代表性。

研究数据来自国家环境监测中心的标准监测网络，包含每日PM$_{2.5}$浓度、站点经纬度坐标等字段~\cite{chang2015outlier}。CMAQ模型输出为空间分辨率36km的格点数据，覆盖研究区域全域~\cite{rodriguez2025downscaling}。

数据预处理包括：（1）剔除缺失值和异常值；（2）将监测站坐标匹配至最近的CMAQ网格点；（3）按十折交叉验证方案划分训练集与测试集。

\subsection{基准方法}

为评估提出方法的性能，本研究引入四种基准融合方法进行比较~\cite{liu2019spatial}：

\begin{itemize}
    \item \textbf{VNA（Voronoi邻居平均）}：直接将监测站观测值插值至CMAQ格点，利用Voronoi图确定最近邻站点~\cite{singh2019enhanced}

    \item \textbf{eVNA（增强型VNA）}：在VNA基础上引入CMAQ作为协变量进行乘法偏差校正

    \item \textbf{aVNA（加法VNA）}：采用加法偏差校正的局地化变分插值

    \item \textbf{Downscaler}：基于贝叶斯推断的MCMC采样降尺度方法~\cite{rodriguez2025downscaling}
\end{itemize}

VNA方法的计算公式为：
\begin{equation}
O_{VNA}(s^*) = \frac{\sum_{i=1}^{k} w_i \cdot O(s_i)}{\sum_{i=1}^{k} w_i}, \quad w_i = d_i^{-2}
\end{equation}
其中$d_i$为预测点$s^*$到第$i$个近邻的距离。

eVNA采用乘法偏差校正：
\begin{equation}
O_{eVNA}(s^*) = O_{CMAQ}(s^*) \times r_n, \quad r_n = \frac{O_{monitor}}{O_{VNA}}
\end{equation}

aVNA采用加法偏差校正：
\begin{equation}
O_{aVNA}(s^*) = O_{CMAQ}(s^*) + bias, \quad bias = O_{monitor} - O_{VNA}
\end{equation}

Downscaler方法利用贝叶斯推断框架，通过MCMC采样联合推断CMAQ网格值与监测站观测值之间的关系~\cite{rodriguez2025downscaling}。

\subsection{AdvancedRK方法}

AdvancedRK方法的整体框架包含三个核心步骤：多项式趋势校正、残差计算与空间建模、以及最终融合预测。该方法的核心创新在于：（1）采用二阶多项式建模CMAQ与监测值之间的非线性偏差关系；（2）利用高斯过程回归（GPR）与Matern核函数对残差进行空间插值。

\subsubsection{Step 1：多项式趋势校正}

设$CMAQ(s_i)$为位置$s_i$处的CMAQ模拟值，$O(s_i)$为对应的监测观测值。本研究采用二阶多项式形式：

\begin{equation}
M_{cal}(s_i) = \alpha_0 + \alpha_1 \cdot CMAQ(s_i) + \alpha_2 \cdot CMAQ(s_i)^2
\end{equation}

多项式系数通过普通最小二乘法（OLS）估计~\cite{rasmussen2006gaussian}：

\begin{equation}
\boldsymbol{\alpha} = (\mathbf{X}^T \mathbf{X})^{-1} \mathbf{X}^T \mathbf{y}
\end{equation}

其中$\mathbf{X}$为设计矩阵，$\mathbf{y}$为观测值向量。

\textbf{为何使用二阶多项式？}

选择二阶多项式基于以下考虑：（1）CMAQ模型与实际监测值之间的偏差往往不是简单的线性关系，而是存在非线性特征~\cite{zheng2020kriging}；（2）二阶多项式能够捕捉U形或抛物线形的偏差关系，即在CMAQ模拟值较低和较高时偏差方向可能不同；（3）更高阶的多项式容易过拟合，且增加计算复杂度。

\subsubsection{Step 2：GPR残差空间建模}

计算每个监测站的残差（即观测值与多项式预测值的差）：

\begin{equation}
R(s_i) = O(s_i) - M_{cal}(s_i)
\end{equation}

对残差序列建立高斯过程先验~\cite{rasmussen2006gaussian}：

\begin{equation}
R(\cdot) \sim \mathcal{GP}(\mu_R, k(\cdot, \cdot'; \theta))
\end{equation}

其中$\mu_R$为残差均值函数，$k$为核函数，$\theta$为核参数。

\textbf{Matern核函数（AdvancedRK采用）：}

\begin{equation}
k_{Matern}(s_i, s_j) = \sigma^2 \frac{1}{\Gamma(\nu)2^{\nu-1}} \left(\frac{\sqrt{2\nu}\|s_i - s_j\|}{\ell}\right)^\nu K_\nu\left(\frac{\sqrt{2\nu}\|s_i - s_j\|}{\ell}\right)
\end{equation}

其中：
\begin{itemize}
    \item $\|s_i - s_j\|$：位置$s_i$和$s_j$之间的欧氏距离
    \item $\ell$：长度尺度参数，控制相关性的作用范围
    \item $\sigma^2$：信号方差，控制输出的整体方差
    \item $\nu$：平滑参数，控制函数的可微性
    \item $K_\nu$：第二类修正贝塞尔函数
\end{itemize}

当$\nu=1.5$时，Matern核等价于：
\begin{equation}
k_{Matern}^{(\nu=1.5)}(s_i, s_j) = \sigma^2 \left(1 + \frac{\sqrt{3}\|s_i - s_j\|}{\ell}\right) \exp\left(-\frac{\sqrt{3}\|s_i - s_j\|}{\ell}\right)
\end{equation}

\textbf{为何选择Matern核（$\nu=1.5$）？}

与传统RBF核相比，Matern核具有以下优势：（1）RBF核假设空间相关性随距离呈指数平方衰减，衰减速度过快，不符合PM$_{2.5}$污染物扩散的实际规律；（2）$\nu=1.5$的Matern核在距离较远时衰减更慢，更符合污染物从源区向外扩散时空间相关性平缓衰减的物理特征；（3）$\nu=1.5$的Matern核函数一阶可微，能够刻画PM$_{2.5}$浓度场的局部变化特征~\cite{rasmussen2006gaussian}。

\textbf{克里金预测公式为：}

给定新位置$s^*$，其残差预测值为~\cite{krige1951note}：

\begin{equation}
R^*(s^*) = \mu_R(s^*) + \mathbf{k}^T \mathbf{K}^{-1} (\mathbf{R} - \mu_R)
\end{equation}

其中：
\begin{itemize}
    \item $\mu_R(s^*)$：新位置$s^*$处的残差均值
    \item $\mathbf{k} = [k_{Matern}(s^*, s_1), k_{Matern}(s^*, s_2), \ldots, k_{Matern}(s^*, s_n)]^T$：新位置与所有训练点之间的协方差向量
    \item $\mathbf{K}$：训练点之间的协方差矩阵，$K_{ij} = k_{Matern}(s_i, s_j)$
    \item $\mathbf{R} = [R(s_1), R(s_2), \ldots, R(s_n)]^T$：训练点的残差向量
\end{itemize}

GPR的核参数（$\ell$和$\sigma$）通过最大化对数边际似然估计确定：

\begin{equation}
\log p(\mathbf{R} | \mathbf{X}) = -\frac{1}{2} \mathbf{R}^T \mathbf{K}^{-1} \mathbf{R} - \frac{1}{2} \log |\mathbf{K}| - \frac{n}{2} \log(2\pi)
\end{equation}

\subsubsection{Step 3：融合预测}

最终融合预测值为多项式校正模型与残差预测之和~\cite{wan2022fusion}：

\begin{equation}
P_{AdvancedRK}(s^*) = M_{cal}(s^*) + R^*(s^*)
\end{equation}

即先通过多项式校正得到CMAQ偏差的全局估计，再通过GPR-Matern克里金补充局部空间信息。

\subsection{性能评估指标}

采用四个标准指标评估融合性能~\cite{liu2019spatial}：

\begin{itemize}
    \item \textbf{决定系数} $R^2$：$1 - \sum(y_i - \hat{y}_i)^2 / \sum(y_i - \bar{y})^2$，衡量模型解释方差的比例，越接近1越好

    \item \textbf{平均绝对误差} MAE：$\frac{1}{n}\sum|y_i - \hat{y}_i|$，对异常值不如RMSE敏感

    \item \textbf{均方根误差} RMSE：$\sqrt{\frac{1}{n}\sum(y_i - \hat{y}_i)^2}$，惩罚大误差

    \item \textbf{平均偏差} MB：$\frac{1}{n}\sum(\hat{y}_i - y_i)$，反映系统性偏差方向
\end{itemize}

\section{实验结果}

\subsection{多阶段验证设置}

本研究采用多阶段验证策略~\cite{rodriguez2025downscaling}，包括：

\begin{itemize}
    \item \textbf{预实验阶段（pre\_exp）}：2020年1月1日至5日（5天），快速筛选方法潜力
    \item \textbf{Stage 1}：2020年1月整月（31天），冬季验证
    \item \textbf{Stage 2}：2020年7月整月（31天），夏季验证
    \item \textbf{Stage 3}：2020年12月整月（31天），冬季验证
\end{itemize}

每天进行十折交叉验证，合并所有天的预测结果计算整体指标。创新判定标准为：$R^2 \geq R^2_{best\_baseline} + 0.01$，且RMSE和|MB|不差于最优基准。

\subsection{基准方法性能对比}

表~\ref{tab:benchmark}展示了各基准方法在多阶段验证中的$R^2$性能表现。

\begin{table}[htbp]
\centering
\caption{基准方法多阶段验证$R^2$对比}
\label{tab:benchmark}
\begin{tabular}{@{}lcccc@{}}
\toprule
方法 & pre\_exp & stage1 & stage2 & stage3 \\
\midrule
VNA & 0.8907 & 0.9034 & 0.8408 & 0.9031 \\
aVNA & 0.8883 & 0.9014 & 0.8175 & 0.9007 \\
eVNA & 0.8842 & 0.8913 & 0.7595 & 0.8924 \\
CMAQ & 0.2682 & 0.2083 & $-$0.3464 & 0.2967 \\
\bottomrule
\end{tabular}
\end{table}

基准方法中，VNA表现最佳，各阶段$R^2$均在0.84--0.90之间。CMAQ原始模拟值效果较差，存在明显低估倾向（MB=$-$12.70 $\mu$g/m$^3$）。eVNA和aVNA方法引入CMAQ作为协变量后并未显著提升性能，表明简单的线性偏差校正假设可能过于简化~\cite{liu2019spatial}。

\subsection{AdvancedRK方法多阶段验证结果}

表~\ref{tab:innovation}展示了AdvancedRK方法在各验证阶段的性能及与VNA基准的对比。

\begin{table}[htbp]
\centering
\caption{AdvancedRK方法多阶段验证结果}
\label{tab:innovation}
\begin{tabular}{@{}lcccccc@{}}
\toprule
阶段 & AdvancedRK & VNA & AdvancedRK & VNA & AdvancedRK & 判定 \\
     & $R^2$      & $R^2$ & RMSE     & RMSE & |MB|       &      \\
\midrule
pre\_exp & 0.9015 & 0.8907 & 15.84 & 16.68 & 0.05 & 通过 \\
stage1   & 0.9143 & 0.9034 & 15.51 & 16.48 & 0.14 & 通过 \\
stage2   & 0.8501 & 0.8408 & 4.90  & 5.05  & 0.04 & 通过 \\
stage3   & 0.9104 & 0.9031 & 11.73 & 12.20 & 0.01 & 通过 \\
\bottomrule
\end{tabular}
\end{table}

AdvancedRK方法在全部四个阶段均通过创新判定。具体而言：
\begin{itemize}
    \item \textbf{预实验阶段}：$R^2$=0.9015，较VNA提升+0.0108，RMSE改善0.84 $\mu$g/m$^3$
    \item \textbf{Stage 1（1月）}：$R^2$=0.9143，较VNA提升+0.0109，RMSE改善0.97 $\mu$g/m$^3$
    \item \textbf{Stage 2（7月）}：$R^2$=0.8501，较VNA提升+0.0093，RMSE改善0.15 $\mu$g/m$^3$
    \item \textbf{Stage 3（12月）}：$R^2$=0.9104，较VNA提升+0.0073，RMSE改善0.47 $\mu$g/m$^3$
\end{itemize}

\subsection{与其他创新方法对比}

表~\ref{tab:compare}展示了AdvancedRK与其他创新方法的对比。

\begin{table}[htbp]
\centering
\caption{创新方法单日验证对比（2020-01-01）}
\label{tab:compare}
\begin{tabular}{@{}lcccc@{}}
\toprule
方法 & $R^2$ & MAE & RMSE & MB \\
\midrule
CMAQ原始值 & $-$0.038 & 20.47 & 29.25 & $-$3.24 \\
VNA & 0.800 & 7.75 & 12.86 & 0.76 \\
eVNA & 0.810 & 7.99 & 12.52 & 0.08 \\
RK-Poly & 0.852 & 7.09 & 11.05 & 0.17 \\
\textbf{AdvancedRK} & \textbf{0.850} & \textbf{7.16} & \textbf{11.11} & \textbf{0.20} \\
\bottomrule
\end{tabular}
\end{table}

\section{讨论}

\subsection{方法分析}

AdvancedRK方法相比VNA基准在多个阶段实现了显著的性能提升。可能原因包括：

\begin{enumerate}
    \item \textbf{非线性建模能力}：二阶多项式趋势校正捕捉了CMAQ与监测值之间的非线性偏差关系~\cite{zheng2020kriging}。CMAQ模型在高浓度和低浓度区域的偏差方向和幅度不同，二阶多项式能够有效刻画这种非线性特征。

    \item \textbf{残差空间结构利用}：GPR能够捕捉残差的空间相关性，对局部偏差进行修正。残差中包含的空间信息通过克里金插值传递至预测点，提高了融合精度。

    \item \textbf{Matern核物理意义}：$\nu=1.5$的Matern核更符合PM$_{2.5}$从源区向外扩散时空间相关性平缓衰减的特征。与RBF核的快速衰减相比，Matern核在距离较远时仍保持一定相关性，符合大气污染物扩散的物理规律。

    \item \textbf{系统性偏差消除}：|MB|在多个阶段接近于零（0.01--0.14 $\mu$g/m$^3$），表明残差建模有效消除了融合结果的系统性偏差。
\end{enumerate}

传统克里金方法在气象和环境数据融合中已有广泛应用~\cite{cressman1959operational,cooley2007kriging}。本研究的结果表明，基于GPR的残差克里金在PM$_{2.5}$数据融合中能够提供有效的精度提升~\cite{zheng2020kriging}。

\subsection{核函数对比分析}

RBF核与Matern核的关键差异在于相关性衰减方式：

\begin{itemize}
    \item \textbf{RBF核}：$k(d) = \exp(-d^2/(2\ell^2))$，距离为$\ell$时相关性衰减至$e^{-0.5}\approx 0.61$，衰减速度过快
    \item \textbf{Matern核($\nu=1.5$)}：$k(d) = (1+\sqrt{3}d/\ell)\exp(-\sqrt{3}d/\ell)$，包含多项式衰减因子，在距离$\ell$时仍有较强相关性
\end{itemize}

对于PM$_{2.5}$污染物扩散问题，Matern核更符合物理规律：污染物从高浓度区向外扩散时，空间相关性不会突然消失，而是随距离平缓衰减。这一特性使得AdvancedRK在空间插值时能够更好地利用远距离站点的信息。

\subsection{方法优势}

AdvancedRK方法展现出以下优势：

\begin{enumerate}
    \item \textbf{多阶段稳定性}：在冬夏不同季节（1月、7月、12月）均保持稳定性能，全部四个阶段通过创新判定

    \item \textbf{非线性建模能力}：二阶多项式趋势校正捕捉了CMAQ与监测值之间的非线性偏差关系

    \item \textbf{系统性偏差消除}：|MB|接近于零表明残差建模有效消除了融合结果的系统性偏差

    \item \textbf{RMSE全面改善}：所有阶段的RMSE均低于VNA基准，最大改善达0.97 $\mu$g/m$^3$

    \item \textbf{物理可解释性}：Matern核的选择基于PM$_{2.5}$空间扩散的物理特征，具有明确的物理解释
\end{enumerate}

\subsection{未来改进方向}

基于本研究的发现，未来可从以下方向改进AdvancedRK方法~\cite{wang2025spint}：

\begin{enumerate}
    \item \textbf{引入时间维度}：扩展至时空融合框架，纳入PM$_{2.5}$的日间时间演变规律

    \item \textbf{多源数据融合}：纳入卫星AOD、气象协变量（温度、湿度、风速）等额外数据源

    \item \textbf{物理约束集成}：将物理规律（如大气扩散方程）作为约束条件引入融合框架

    \item \textbf{自适应核选择}：根据局部数据特征自动选择最优核函数参数
\end{enumerate}

\section{结论}

本研究针对PM$_{2.5}$ CMAQ数据融合问题，提出了基于多项式残差克里金与Matern核高斯过程回归的AdvancedRK融合方法。该方法的核心贡献包括：

\begin{enumerate}
    \item 采用二阶多项式建模CMAQ与监测值之间的非线性偏差关系，提高了趋势校正的准确性

    \item 利用GPR-Matern克里金对残差进行空间插值，捕捉残差的复杂空间相关性

    \item 引入Matern核函数（$\nu=1.5$），更符合PM$_{2.5}$空间扩散的物理特征，优于传统RBF核

    \item 通过多阶段十折交叉验证严格评估方法性能，确保结果的统计可靠性
\end{enumerate}

基于中国某地区1528个监测站的多阶段验证表明，AdvancedRK方法在全部四个阶段通过创新判定：pre\_exp阶段$R^2$=0.9015（VNA：0.8907），stage1阶段$R^2$=0.9143（VNA：0.9034），stage2阶段$R^2$=0.8501（VNA：0.8408），stage3阶段$R^2$=0.9104（VNA：0.9031）。综合四阶段验证结果，AdvancedRK方法显著超越所有基准方法，验证了其有效性和创新性~\cite{rodriguez2025downscaling}。

本研究的发现表明，多项式残差克里金方法结合Matern核函数对PM$_{2.5}$融合精度的提升有显著贡献。未来研究应考虑引入时间维度、多源卫星数据或物理约束，以进一步提升融合方法的性能。

\section*{致谢}

本研究使用了国家环境监测中心提供的PM$_{2.5}$监测数据，以及CMAQ模型模拟数据。感谢各位审稿专家提出的宝贵修改意见。

\bibliographystyle{plain}
\bibliography{references}

\end{document}
